\chapter{Native Extensions}\label{ch:pl:native-extensions}

A researcher rewrote the inner step of a signal in C++ and called it from Python once per tick, a million times. It was slower
than the Python function it replaced: 118 nanoseconds a call against 60. Called once on the whole array of a million values, the
same C++ loop took 1.4 nanoseconds an element, more than twice as fast as the library routine it was meant to beat. Everything
about a \omterm{def:pl:native-extensions:extension}{native extension} is decided at the boundary between the interpreter and the compiled code: how often it is crossed, what
has to be converted on the way, and who owns the memory on each side. This chapter binds one kernel written in C++20 and in Rust into
Python two ways, with pybind11 and through a C interface loaded with \texttt{ctypes}, passes numpy arrays across without copying them,
and measures the boundary (\texttt{firm.natext}).

\section{Why and when to leave Python}

\begin{definition}[Native extension, language binding]\label{def:pl:native-extensions:extension}
A \emph{native extension}\index{native extension} is compiled code (C, C++, Rust) loaded into the Python process and called like a
Python function. A \emph{language binding}\index{language binding} is the layer that makes it callable: it converts arguments and
results between the two languages' representations, reports errors across the boundary, and states who owns each piece of memory.
\end{definition}

Chapter~8 gave the order of preference: first \omterm{def:pl:python-at-scale:array}{array programming}, which is already compiled code written by someone else; then a
\omterm{def:pl:native-extensions:extension}{native extension} for what arrays cannot express --- a loop whose state changes with each element (a book, a queue position, a state
machine), a merge of two sorted sequences, a recursion with branches. The chapter's two kernels are one of each kind: an exponentially
weighted average (which numpy can express as a filter, so the native version competes with a library) and an as-of lookup of sorted
times by one merge pass (which numpy does with a binary search per element, $O(n \log m)$ instead of $O(n + m)$).

\omcode{code/firm/natext/cpp/firm_natext.hpp}{9}{38}{The two kernels in C++20, free of Python: the average, one step of it, and the
as-of lookup by one merge pass.}\label{lst:pl:native-extensions:kernels}

\section{The boundary: calls, types and the interpreter lock}

\begin{definition}[Foreign function interface, application binary interface]\label{def:pl:native-extensions:ffi}
A \emph{foreign function interface}\index{foreign function interface} lets code in one language call functions compiled from another.
It relies on an \emph{application binary interface}\index{application binary interface}: the machine-level convention (how arguments
are passed in registers and on the stack, the layout of structures, the names of symbols) that the caller and the callee compiled
separately both obey; the C convention is the one every language can speak.
\end{definition}

\begin{definition}[Boundary-crossing cost]\label{def:pl:native-extensions:crossing}
The \emph{boundary-crossing cost}\index{boundary-crossing cost} of a native call is the fixed work done on each call, independently of
the size of its data: finding the function, checking and converting each argument from a Python object, releasing and reacquiring the
interpreter lock, converting the result back and handling errors.
\end{definition}

\begin{omfigure}
\begin{tikzpicture}[font=\scriptsize,>=Stealth,box/.style={draw,rounded corners=2pt,minimum height=0.7cm,inner sep=3pt,align=center}]
\node[box,fill=omDef!12] (py) at (0,0) {Python caller:\\objects, numpy arrays};
\node[box,fill=omMeth!15] (pb) at (4.6,0.9) {pybind11 binding: checks\\type and layout, releases\\the lock};
\node[box,fill=omThm!12] (ct) at (4.6,-0.9) {ctypes: passes pointers\\as declared, checks nothing\\(the wrapper must)};
\node[box,fill=gray!10] (cpp) at (9.0,0.9) {C++20 kernel};
\node[box,fill=gray!10] (rs) at (9.0,-0.9) {Rust kernel\\(C ABI, cdylib)};
\draw[->] (py) -- (pb); \draw[->] (py) -- (ct); \draw[->] (pb) -- (cpp); \draw[->] (ct) -- (rs);
\draw[dashed,omThm,rounded corners=3pt] (-1.6,-2.4) rectangle (10.4,-1.8);
\node[omThm] at (4.4,-2.1) {one numpy buffer: read and written in place by both kernels, never copied};
\end{tikzpicture}

\omcaption{Two routes across the boundary. pybind11 generates a binding that checks each argument's type and memory layout and can
release the interpreter lock while C++ runs; ctypes loads any library with a C interface and passes what it is told, so the Python
wrapper must check. Both hand the kernel a pointer into the same numpy buffer.}\label{fig:pl:native-extensions:boundary}
\end{omfigure}

\Cref{fig:pl:native-extensions:calls} measures the crossing. A one-step function called a million times from a Python loop costs 60
nanoseconds a call when it is a Python function, 118 through pybind11 and 337 through \texttt{ctypes}: each native call converts
three Python floats to C doubles and one back, and \texttt{ctypes} does it dynamically, by the argument types declared at run time.
The arithmetic, one multiplication and one addition, is less than a nanosecond. A native function called per element is slower than
the Python it replaces, because the loop around it is still Python.

\begin{omfigure}
\begin{tikzpicture}
\begin{axis}[width=11cm,height=4.4cm,xbar,bar width=9pt,xmin=0,xmax=480,xlabel={nanoseconds per call, loop included},y dir=reverse,
  ytick={0,1,2},yticklabels from table={figdata/platforms/09-native-extensions/measured_calls.csv}{route},
  tick label style={font=\small},label style={font=\small},grid=major,grid style={gray!20},point meta=explicit,
  nodes near coords,nodes near coords style={font=\footnotesize,anchor=west,/pgf/number format/fixed,/pgf/number format/precision=0},
  table/col sep=comma]
\addplot[fill=omThm!45,draw=omThm] table[x=ns_per_call,y expr=\coordindex,meta=ns_per_call]{figdata/platforms/09-native-extensions/measured_calls.csv};
\end{axis}
\end{tikzpicture}

\omcaption{The cost of calling a one-step function once per element from a Python loop, a million times: as a Python function, as C++
through pybind11, and as Rust through \texttt{ctypes}. Measured on a laptop (Intel Core Ultra 7 155H) under WSL2, one thread, machine
otherwise idle. Data: \texttt{bench\_natext.py}.}\label{fig:pl:native-extensions:calls}
\end{omfigure}

\begin{proposition}[Move the loop, not the body]\label{prop:pl:native-extensions:loop}
If a native call costs $c_0$ per crossing and $c_1$ per element, and the Python it replaces costs $p$ per element, then calling it
once per element costs $n(c_0 + c_1)$ and is faster than Python only if $c_0 + c_1 < p$; calling it once on the whole array costs
$c_0 + n c_1$, which is faster than Python for every $n > c_0/(p - c_1)$ and approaches $n c_1$.
\end{proposition}

\begin{proof}
Sum the costs of the calls in each arrangement and compare with $np$.
\end{proof}

With $c_0 \approx 118$~ns and $p \approx 60$~ns, the first arrangement never wins; the second wins beyond a handful of elements. The
interpreter lock is the other half of the boundary: while the kernel runs it touches no Python object, so the binding releases the
lock (\cref{lst:pl:native-extensions:pybind}) and other Python threads run meanwhile (chapter~8's sort did the same).

\section{Binding C++ with a header-only binder}

pybind11 is a header-only C++ library: the binding is a C++ file that includes it, describes each function's arguments, and is compiled
with the Python headers into a shared library that Python imports as a module. \texttt{firm.natext} compiles it on first use with the
compiler of the series (\texttt{g++ -std=c++20 -O2 -shared -fPIC} with the include paths from \texttt{python3-config} and
pybind11).

\omcode{code/firm/natext/cpp/firm_natext_py.cpp}{11}{24}{The pybind11 binding of the average: a numpy array taken as a C-contiguous
float64 buffer, a result array allocated once, and the interpreter lock released while C++ runs.}\label{lst:pl:native-extensions:pybind}

\begin{definition}[Buffer protocol]\label{def:pl:native-extensions:buffer}
The \emph{buffer protocol}\index{buffer protocol} is Python's interface by which an object exposes its underlying memory --- a pointer,
the element type, the shape and the strides --- to other objects and to compiled code, so that they can read and write it without
copying.
\end{definition}

A numpy array exposes its memory this way, and pybind11's \texttt{array\_t} reads it. The declaration decides what happens when the
array is not what the kernel expects. With \texttt{forcecast}, pybind11 silently converts: a float32 array or a strided view is copied
into a new contiguous float64 buffer, and a function that was meant to write into the caller's array writes into the copy instead. The
chapter's binding declares \texttt{c\_style} without \texttt{forcecast} and marks its arguments \texttt{noconvert}, so both are refused
with an error: a copy should be a decision, not a surprise.

\section{Binding Rust through the C ABI}

Rust has a binding generator for Python too; this series uses no external crates, so the chapter takes the route that needs none: a
library of \texttt{extern \textquotedbl{}C\textquotedbl{}} functions compiled as a C-compatible dynamic library, loaded with Python's \texttt{ctypes}.

\omcode{code/firm/natext/rust/src/lib.rs}{29}{40}{A Rust function with a C interface: raw pointers and a length in, a slice built from
them under stated conditions, and the safe kernel called on the slices.}\label{lst:pl:native-extensions:rust}

\begin{dated}{2026-09}{Rust's Python bindings}\label{dat:pl:native-extensions:pyo3}
PyO3, described in its user guide as ``Rust bindings for Python, including tools for creating native Python extension modules''
(consulted September 2026), generates the equivalent of pybind11's argument checks and conversions for Rust. It is an external crate,
which the series does not build; the C interface of this chapter is what such a binding produces underneath. A \texttt{cdylib} is, in
the Rust Reference's words, a dynamic system library used when compiling a library to be loaded from another language.
\end{dated}

\texttt{ctypes} is the other extreme from pybind11: it calls any C function with the arguments it is told to pass and checks nothing
about the memory behind a pointer. The wrapper (\cref{lst:pl:native-extensions:ctypes}) declares each function's argument types once,
checks that each array is C-contiguous and of the right type, and passes the pointers; without the check, a strided view would be read
as if it were contiguous, silently.

\omcode{code/firm/natext/firm_natext.py}{105}{131}{Loading the Rust library with \texttt{ctypes}: declared argument types, and a
wrapper that checks what \texttt{ctypes} does not before passing pointers into numpy buffers.}\label{lst:pl:native-extensions:ctypes}

\section{Zero-copy across the boundary}

\begin{definition}[Zero-copy transfer]\label{def:pl:native-extensions:zerocopy}
A \emph{zero-copy transfer} of data across a boundary\index{zero-copy transfer} hands the receiver a reference to the sender's memory
(a pointer and a layout) instead of a copy, so that its cost does not depend on the size of the data; it requires the two sides to agree
on the layout and on who may write and free the memory, and for how long.
\end{definition}

Both bindings are zero-copy for their inputs; \texttt{ewma\_into} is zero-copy for its output as well, writing into an array the
caller allocated (a test checks that the array's address is unchanged). The ownership rules are what makes it safe: the kernel reads
and writes only during the call, keeps no pointer afterwards, and never frees what Python allocated; Python keeps the arrays alive for
the call's duration. Chapter~4's \omterm{def:pl:a-tick-store-on-open-formats:mmap}{memory-mapped files} extend the same idea to disk and to other processes, and Arrow's layout of
chapter~3 to other libraries and languages.

\begin{omfigure}
\begin{tikzpicture}
\begin{axis}[width=12cm,height=6cm,xmode=log,ymode=log,xlabel={array length (log scale)},ylabel={ns per element (log scale)},
  tick label style={font=\small},label style={font=\small},grid=major,grid style={gray!20},
  legend columns=4,legend style={font=\footnotesize,at={(0.5,-0.25)},anchor=north,draw=none,/tikz/every even column/.append style={column sep=6pt}},
  table/col sep=comma]
\addplot[thick,mark=*,omDef] table[x=n,y=ewma_numpy]{figdata/platforms/09-native-extensions/measured_sweep.csv};
\addplot[thick,mark=square*,omMeth] table[x=n,y=ewma_cpp]{figdata/platforms/09-native-extensions/measured_sweep.csv};
\addplot[thick,mark=triangle*,omThm] table[x=n,y=ewma_rust]{figdata/platforms/09-native-extensions/measured_sweep.csv};
\addplot[thick,mark=diamond*,gray] table[x=n,y=ewma_python]{figdata/platforms/09-native-extensions/measured_sweep.csv};
\legend{numpy (filter),C++ (pybind11),Rust (ctypes),pure Python}
\end{axis}
\end{tikzpicture}

\omcaption{Time per element of the whole-array average against the array's length. For short arrays every route is dominated by its
fixed cost per call ($c_0/n$); for long arrays the native loops settle at about 1.3 to 1.9~ns an element and numpy's filter at 3 to 3.6, while
pure Python stays at hundreds. Measured on a laptop (Intel Core Ultra 7 155H) under WSL2, one thread, machine otherwise idle. Data:
\texttt{bench\_natext.py}.}\label{fig:pl:native-extensions:sweep}
\end{omfigure}

\Cref{fig:pl:native-extensions:sweep,fig:pl:native-extensions:asof} are the whole-array versions. For the average, both native routes
beat numpy's filter at every length measured: the filter routine has a fixed cost of several microseconds, larger than pybind11's; at a
million elements C++ takes 1.4 nanoseconds an element against numpy's 3.3. For the as-of lookup the algorithm matters more than the
language: numpy's binary search per element costs more per element as the arrays grow, 21~ns at a million, while the merge pass stays at
about 6, three times less. The \texttt{ctypes} route pays a larger fixed cost (the pointer conversions of the wrapper) and overtakes
numpy only from about a thousand elements for the lookup.

\begin{omfigure}
\begin{tikzpicture}
\begin{axis}[width=12cm,height=5.6cm,xmode=log,ymode=log,xlabel={array length (log scale)},ylabel={ns per element (log scale)},
  tick label style={font=\small},label style={font=\small},grid=major,grid style={gray!20},
  legend columns=3,legend style={font=\footnotesize,at={(0.5,-0.25)},anchor=north,draw=none,/tikz/every even column/.append style={column sep=6pt}},
  table/col sep=comma]
\addplot[thick,mark=*,omDef] table[x=n,y=asof_numpy]{figdata/platforms/09-native-extensions/measured_sweep.csv};
\addplot[thick,mark=square*,omMeth] table[x=n,y=asof_cpp]{figdata/platforms/09-native-extensions/measured_sweep.csv};
\addplot[thick,mark=triangle*,omThm] table[x=n,y=asof_rust]{figdata/platforms/09-native-extensions/measured_sweep.csv};
\legend{numpy (binary search per element),C++ (merge),Rust (merge)}
\end{axis}
\end{tikzpicture}

\omcaption{Time per element of the as-of lookup of $n$ sorted times in $n$ others. numpy's binary search costs more per element as the
arrays grow; the native merge pass costs a constant few nanoseconds once the fixed cost is amortised. Measured on a laptop (Intel Core
Ultra 7 155H) under WSL2, one thread, machine otherwise idle. Data: \texttt{bench\_natext.py}.}\label{fig:pl:native-extensions:asof}
\end{omfigure}

\begin{method}[Writing a native kernel for the research platform]\label{met:pl:native-extensions:method}
(1) Write the kernel in C++ or Rust with no knowledge of Python: pointers, lengths, plain types. (2) Test it alone in its own language.
(3) Bind it once per array, never once per element: arrays in, arrays out, the loop inside. (4) Refuse layouts and types the kernel does
not read in place, rather than copying silently. (5) Release the interpreter lock around the pure computation. (6) Test the binding
against the Python reference on the same data, including empty and edge inputs.
\end{method}

\section{Tutorial: one kernel, four implementations}\label{tut:pl:native-extensions:tut}

\begin{tutorial}
\textbf{Goal.} Build the C++ and Rust kernels, bind them, show that four implementations agree, and measure the boundary.
\textbf{End state:} \cref{fig:pl:native-extensions:calls,fig:pl:native-extensions:sweep,fig:pl:native-extensions:asof}.
\begin{tutsteps}
\item \textbf{Kernels}: read \cref{lst:pl:native-extensions:kernels}; compile and run \texttt{firm\_natext\_test.cpp}; run
\texttt{cargo test} in \texttt{code/firm/natext/rust}.
\item \textbf{Bind}: \texttt{firm\_natext.cpp\_module()} compiles and imports the pybind11 module;
\texttt{rust\_library()} builds the \texttt{cdylib} and declares its functions to \texttt{ctypes}.
\item \textbf{Agree}: the acceptance tests compare the four averages and the four lookups, on random and edge inputs.
\item \textbf{Refuse}: pass a strided view and a float32 array to each binding and read the errors.
\item \textbf{Measure}: \texttt{bench\_natext.py} (per-call cost and the size sweep).
\end{tutsteps}
\textbf{What to change next.} Add \texttt{forcecast} to the binding and watch \texttt{ewma\_into} write into a copy; call the Rust kernel
from two Python threads through \texttt{ctypes} (which releases the lock around foreign calls) and measure the speed-up.
\end{tutorial}

\section{Build: native kernels for the research platform}\label{bld:pl:native-extensions:natext}

\begin{build}
\bfield{Purpose} The pattern and the build machinery every native kernel of the platform follows: the \omterm{def:pl:a-tick-store-on-open-formats:store}{tick store}'s readers of
chapter~4, the backtest engine's inner loops (chapter~11) and the payoff compiler's evaluation (chapter~19) can all be moved across the
boundary this way.
\bfield{Interface} \texttt{ewma\_python}, \texttt{ewma\_numpy}, \texttt{ewma\_cpp}, \texttt{ewma\_cpp\_into}, \texttt{ewma\_rust};
\texttt{asof\_python}, \texttt{asof\_numpy}, \texttt{asof\_cpp}, \texttt{asof\_rust}; \texttt{step\_python}, \texttt{step\_cpp},
\texttt{step\_rust}; \texttt{cpp\_module()}, \texttt{rust\_library()}. C++20 header \texttt{firm::natext} (\texttt{ewma},
\texttt{ewma\_step}, \texttt{asof\_index}); Rust crate \texttt{natext\_rs} with \texttt{natext\_ewma}, \texttt{natext\_asof},
\texttt{natext\_ewma\_step}.
\bfield{Rules} Kernels know nothing of Python; one call per array; inputs read in place or refused; the lock released around the
computation; no pointer kept after a call; builds are part of the tests and a missing toolchain fails them.
\bfield{Acceptance tests} \texttt{code/firm/natext/tests/}, \texttt{cpp/firm\_natext\_test.cpp}, \texttt{rust/}: four averages and four
lookups equal, edges included; three steps equal; output written in place; strided and wrongly typed arrays refused by both bindings.
\bfield{Stretch} A kernel over the \omterm{def:pl:a-tick-store-on-open-formats:store}{tick store}'s flat records (a structured buffer, chapter~4); a second-order filter in all four
languages; the bindings compared under two threads.
\end{build}

\begin{omsources}
\item pybind11 documentation, \emph{NumPy} (buffer protocol, \texttt{array\_t}, \texttt{c\_style}, \texttt{forcecast},
\texttt{noconvert}).
\item T.~Oliphant and C.~Banks, \emph{PEP 3118 -- Revising the buffer protocol} (Final).
\item The Rust Reference, \emph{Linkage} (\texttt{cdylib}); PyO3 user guide.
\end{omsources}

\section{Exercises}

\begin{exercise}[$\star$]\label{exo:pl:native-extensions:1}
At 118~ns a call through pybind11 and 60~ns a step in Python, how long does a day of two hundred million steps take each way, called
per element? And computed by one call at 1.4~ns an element?
\end{exercise}

\begin{exercise}[$\star$]\label{exo:pl:native-extensions:2}
Why is the \texttt{ctypes} call slower than the pybind11 call, although both end in the same machine instructions?
\end{exercise}

\begin{exercise}[$\star$]\label{exo:pl:native-extensions:3}
What does \texttt{forcecast} do with a float32 array passed to a function declared for float64, and why is that dangerous for
\texttt{ewma\_into}?
\end{exercise}

\begin{exercise}[$\star\star$]\label{exo:pl:native-extensions:4}
With \cref{prop:pl:native-extensions:loop}, $c_0 = 118$~ns, $c_1 = 1.4$~ns and $p = 60$~ns, from how many elements is one call on the
whole array faster than the Python loop?
\end{exercise}

\begin{exercise}[$\star\star$]\label{exo:pl:native-extensions:5}
Why does numpy's as-of lookup cost more per element as the arrays grow, and the merge pass not?
\end{exercise}

\begin{exercise}[$\star\star$]\label{exo:pl:native-extensions:6}
A kernel keeps a pointer to a numpy array between calls to reuse it. Describe how that goes wrong, and what the ownership rule forbids.
\end{exercise}

\begin{exercise}[$\star\star\star$]\label{exo:pl:native-extensions:7}
\emph{Coding.} With the chapter's functions, compute the as-of lookup of 100\,000 random sorted times in 50\,000 others by all four
implementations, and check that they agree, including a left time before the first right time and one after the last.
\end{exercise}

\begin{exercise}[$\star\star\star$]\label{exo:pl:native-extensions:8}
\emph{Find the flaw.} ``Our C++ extension takes a list of Python floats, converts it to a \texttt{std::vector<double>}, computes, and
returns a new list. It is faster than numpy on our benchmark of ten numbers, so we use it everywhere.''
\end{exercise}

\section{Problem: The Million Calls}

\begin{problem}[{Weekend problem --- where to cross the boundary}]\label{pb:pl:native-extensions:1}
The chapter's kernels, their four implementations, the two bindings and the measurements of \texttt{bench\_natext.py}.

\medskip\noindent\textbf{Part I --- The crossing.}
\begin{enumerate}
\item What does a native call do besides the arithmetic?
\item What are the measured costs per call of the three routes?
\item Why is a per-element native call slower than the Python function it replaces?
\item What does releasing the interpreter lock allow, and when is it safe?
\item What is an \omterm{def:pl:native-extensions:ffi}{application binary interface}, and why is the C one the common ground?
\end{enumerate}

\medskip\noindent\textbf{Part II --- The two bindings.}
\begin{enumerate}[resume]
\item What does pybind11 check that \texttt{ctypes} does not?
\item What must the \texttt{ctypes} wrapper check itself?
\item What happens to a strided view passed to each binding?
\item Why is \texttt{forcecast} not used?
\item What would PyO3 add to the Rust route?
\end{enumerate}

\medskip\noindent\textbf{Part III --- Arrays.}
\begin{enumerate}[resume]
\item What are the costs per element of the whole-array average at a million elements, by route?
\item Why do the native routes beat numpy's filter at every length measured?
\item What are the costs per element of the as-of lookup at a million, by route?
\item From what length does the Rust lookup beat numpy's?
\item Which of the two kernels gains more from going native, and why?
\end{enumerate}

\medskip\noindent\textbf{Part IV --- The verdict.}
\begin{enumerate}[resume]
\item State the \emph{named result}: the per-call cost of each binding, the lengths above which each native version beats numpy, and the
saving of moving the loop, not the body, across the boundary.
\item How much time does a day of two hundred million steps save by moving the loop into C++?
\item When is a \omterm{def:pl:native-extensions:extension}{native extension} not worth writing?
\item What does the ownership rule require of the kernel and of the caller?
\item In one sentence: where should the boundary between Python and native code be?
\end{enumerate}
\end{problem}

\section{Interview questions}

\begin{interviewq}[$\star$ \iqroles{developer, researcher}]\label{iq:pl:native-extensions:1}
You wrote a function in C++ to speed up a Python loop, and it got slower. What happened?
\end{interviewq}

\begin{interviewq}[$\star\star$ \iqroles{developer}]\label{iq:pl:native-extensions:2}
How do you pass a numpy array to C++ or Rust without copying it, and what can go wrong?
\end{interviewq}

\begin{interviewq}[$\star\star$ \iqroles{developer}]\label{iq:pl:native-extensions:3}
Compare pybind11, ctypes and a Rust binding generator for exposing a numerical kernel to Python.
\end{interviewq}

\begin{interviewq}[$\star\star$ \iqroles{developer}]\label{iq:pl:native-extensions:4}
Why should a \omterm{def:pl:native-extensions:extension}{native extension} release the interpreter lock, and what must it not do while it is released?
\end{interviewq}

\begin{interviewq}[$\star\star\star$ \iqroles{developer, researcher}]\label{iq:pl:native-extensions:5}
Design the native layer of a research platform: which kernels go native, how they are bound, built, tested and distributed.
\end{interviewq}
